\section{当前价格对外部 EPS 的机械交叉核验}

截至 2026 年 7 月 22 日，雅化收盘价为 16.79 元，总股本 11.5256 亿股，对应市值 193.52 亿元。将该价格除以三份完整券商原始报告的 2026E EPS 2.01--2.63 元，得到 6.38--8.35 倍的机械前瞻 PE 区间。这一计算只描述当前价格相对外部预期的位置；它既没有选择合理倍数，也没有验证 EPS 分母，因此不产生 AStock fair value、目标价、上行空间或评级。

\begin{exhibitbox}[当前价格与外部 EPS 的交叉核验：非自有估值]
\begin{tabularx}{\exhibitboxwidth}{L{2.4cm}R{2.5cm}R{2.6cm}R{2.4cm}X}
\toprule
外部机构 & 2026E 归母净利 & 2026E EPS & 16.79 元对应 PE & 关键限制 \\
\midrule
国信证券 & 23.15 亿元 & 2.01 元 & 8.35 倍 & 4 月报告，H2 假设尚未验证 \\
开源证券 & 25.69 亿元 & 2.23 元 & 7.53 倍 & 4 月报告，非公司指引 \\
东吴证券 & 30.32 亿元 & 2.63 元 & 6.38 倍 & 7 月报告，H2 要求最激进 \\
\bottomrule
\end{tabularx}
\end{exhibitbox}
\sourcenote{国信、开源、东吴原始报告，日期和页码见第 2 章；市场价格、股本和市值见东财快照（2026-07-23 抓取）。PE 为 16.79 元除以外部 EPS 的算术结果。}

\section{为什么“6--8 倍”不能直接等于便宜}

低于外部 EPS 所对应的倍数可以有多种含义：市场可能相信 H1 是峰值或库存收益，可能对锂价、进口和行业库存更谨慎，也可能要求现金流和资源成本得到验证。反过来，它也可能低估了 H2 持续性。三种解释都无法仅通过一个 PE 数字区分；正式 H1 的收入、现金流、应收、存货、分部毛利和锂盐单位经济才是判别器。

\begin{sourcequalitybox}[估值边界]
本报告不设置熊/基/牛目标价，也不把现价复制为目标价。完整单票估值至少需要正式 H1 三表、锂盐量--价--成本到利润的桥、民爆收入/毛利/现金桥、税费和少数股东口径，以及可比公司估值的统一口径。以上条件在本报告截止日均未齐备。
\end{sourcequalitybox}

\section{现价的三种叙事，正式 H1 将排除其中至少一种}

\begin{tabularx}{\linewidth}{L{3.2cm}X X}
\toprule
市场叙事 & 16.79 元的低 PE 可能意味着什么 & 正式 H1 的判别证据 \\
\midrule
周期高点折价 & 市场不相信 Q2 锂价、库存或价差可持续 & H1 实际 ASP、毛利、库存和 H2 订单/销量信息 \\
现金流质量折价 & 市场认为利润需要过多应收、存货或合同资产 & H1 经营现金流与营运资本的变动解释 \\
盈利持续性被低估 & 市场低估了资源、成本与民爆对 H2 的支撑 & 分部利润、实际供给/成本、民爆收入与回款 \\
\bottomrule
\end{tabularx}

以上是竞争性解释，不是对资金行为的断言。它们共同说明：外部 EPS 与当前 PE 的算术关系无法替代对 H1 质量的验证。只有当正式披露排除了“周期高点”或“现金流折价”解释，低 PE 才可能具有明确的基本面含义。

\section{转入全年模型前的输入合同}

这张表不是预测参数清单，而是从事件研究升级为完整盈利模型的最低信息门槛。任何一项长期缺失，研究结论都必须停留在“利润弹性已发生、持续性未证实”，不能用外部 EPS 或低 PE 填补空白。

\begin{tabularx}{\linewidth}{L{3.0cm}X X}
\toprule
输入 & 必需原因 & 没有它时的处理 \\
\midrule
锂盐销量、产品结构、实际 ASP & 将行业价格背景转换为公司收入 & 不预测 H2 收入 \\
原料/加工成本、套保和减值 & 判断 Q2 利润是否可持续 & 不给成本优势信用 \\
资源实际交付、自供比例和现金成本 & 判断供应保障是否改变波动率 & 只作为供应框架描述 \\
民爆收入/毛利、订单与回款 & 验证下行保护与现金贡献 & 不上调民爆盈利 \\
经营现金流、应收、存货、合同资产 & 判断利润质量与营运资本占用 & 不发布目标价 \\
\bottomrule
\end{tabularx}

正式 H1 发布后，应先用这五类字段重建 H1 的量--价--成本--现金桥，再判断外部全年利润是否仍具有可实现性；顺序不能倒置。
